%%%%%%%%%%%%%%%%%%%%%%%%%%%%%%%%%%%%%%%%%%%%%%%%%%%%%%%%%%%%%%%%%%%%%%
% How to use writeLaTeX: 
%
% You edit the source code here on the left, and the preview on the
% right shows you the result within a few seconds.
%
% Bookmark this page and share the URL with your co-authors. They can
% edit at the same time!
%
% You can upload figures, bibliographies, custom classes and
% styles using the files menu.
%
%%%%%%%%%%%%%%%%%%%%%%%%%%%%%%%%%%%%%%%%%%%%%%%%%%%%%%%%%%%%%%%%%%%%%%

\documentclass[12pt]{article}

\usepackage{sbc-template}

\usepackage{graphicx,url}

%\usepackage[brazil]{babel}   
\usepackage[utf8]{inputenc}  

     
\sloppy

\title{Exercício 2 Sistemas Distribuidos}

\author{João Roberto Lemos Fidellis\inst{1}, Christian Aguiar Plentz\inst{2} }


\address{Ciência da Computação - Universidade do Vale do Rio dos Sinos (UNISINOS)\\
  Caixa Postal 275 -- 93.022-750 -- São Leopoldo -- RS -- Brazil
\email{\{plentzchristian, jrfidellis\}@edu.unisinos.br}
}

\begin{document} 

\maketitle

\section{Dê exemplos de sistemas (implementações) onde posso ou não ter
acoplamento de tempo e de  espaço, Aqui eu tenho 4 possibilidades (8 exemplos ao total)}
\label{sec:exercicio1}

\subsection{Acoplamento de Tempo, Acoplamento de Espaço}
\textbf{Comunicação entre Sockets TCP/IP} 
Se precisa saber o Endereço IP e a Porta Destinatária do remetente e os
    2 lados precisam estar online para ocorrer a comunicação.
\textbf{Remote Procedure Call (RPC)} 
Se chama uma função remotamente em um servidor que precisa retornar
    essa função para o chamador, é um processo bloqueante que precisa
    do Endereço do servidor e os 2 lados precisam estar ativos durante
    essa operação.

\subsection{Desacoplamento de Tempo, Acoplamento de Espaço}
\textbf{Serviço de E-Mail:} 
O mandante da mensagem precisa saber o endereço de E-Mail do remetente mas os 2
não precisam estar online quando a mensagem for enviada a um intermediario
que armazena os E-Mails 
\textit{Exemplos: Outlook e Gmail}


\noindent \textbf{Serviço de Chat Ponto a Ponto} 
O mandante da mensagem só precisa saber o usuário do remetente, mas não precisa os 2 ficarem
online na mesma hora para o remetente receber a mensagem
já que elas ficam armazenadas no servidor.

\subsection{Acoplamento de Tempo, Desacoplamento de Espaço}
\textbf{IRC (Internet Relay Chat)}
O mandante da Mensagem precisa somente estar na mesma sala com o remetente mas
os 2 precisam estar online ao mesmo tempo para receber a mensagem

\noindent \textbf{Serviço de Reunião Online}
Os integrantes precisam estar online em uma reunião ao mesmo tempo mas não precisam
saber o endereço dos integrantes.
\textit{Ex: Microsoft Teams e Google Meet}
\subsection{Desacoplamento de Tempo, Desacoplamento de Espaço}

\textbf{Espaço de Tuplas} Se tem uma área onde pode ser colocado Objetos Java onde
integrantes podem ler e escrever sobre para recuperar o Objeto ou adicionar novos
objetos. Não se tem necessidade de saber o endereço de quem esta manuseando os dados
e os mesmos ficam salvos dentro do espaço. \textit{Ex: JavaSpaces}

\noindent \textbf{PubSub}
Se tem um servidor com tópicos onde se pode se inscrever e publicar, os produtores
publicam alguma informação no tópico para quando o consumidor se inscrever neste
tópico poder ler as informções publicadas, \textit{Ex: MQTT e RSS}

\section{
Escreva o que você entende pelo tipo de middleware e procure 3 players usados no mercado hoje em dia para cada um deles}

\subsection{PubSub}
Um broker recebe eventos de instrumentos externos para tópicos onde pode então
notificar clientes assinantes ou armazenar para notificar futuros assinantes.
\textit{Players: Google Cloud Pub/Sub, Kafka e RSS}

\subsection{Message Queues}
Um broker recebe mensagens para por em uma fila especifica onde produtores
e consumidores interagem entre si.
\textit{Players: RabbitMQ, ZeroMQ e Apache Qpid}

\subsection{Distributed Shared Memory}
Computadores em cluster que agem como uma unidade de memoria única para serem
utilizados por processos.
\textit{Players: Apache Ignite, Hazelcast e Infinispan}

\subsection{Espaco de Tuplas}
Um espaço onde clientes podem escrever Objetos Java sobre e Ler também
\textit{Players: GigaSpaces e Apache River}

\end{document}

